% II. Анализ на възможните решения и обосновка на избора.
% Тук се сравняват трите нива на достъп и се обосновава защо проектът избира
% най-ниското. Обосновката трябва да назове и цената на избора, не само ползата.
\section{Анализ на възможните решения и обосновка на избора}
\label{sec:alternatives}

\subsection{Разглеждани алтернативи}

Разгледани са три решения на една и съща задача, а именно да се изпълни заявление от
програма на C++ и да се получат резултатните редове в типизиран вид.

\textbf{Алтернатива А, слоест диспечер на драйвери.} Приложението извиква ODBC, а
диспечерът зарежда драйвер за конкретния сървър~\cite{odbc_ref}. Предимството е една
входна точка за много източници на данни, включително такива, за които никой не би
написал клиент сам. Недостатъкът е, че всеки резултатен ред минава през договорения от
стандарта тип и през буфер, зададен от приложението предварително, а асинхронната
работа не е част от модела.

\textbf{Алтернатива Б, доставчикова клиентска библиотека.} Приложението се свързва
пряко с \code{libpq} или с \code{libmysqlclient}. Отпада един слой, а семантиката е
точно тази на сървъра. Остават обаче два въпроса: библиотеките са на C и налагат
собствено управление на паметта на границата, а асинхронните им режими се различават
помежду си достатъчно, за да не се обединяват просто~\cite{libpq}.

\textbf{Алтернатива В, пряка реализация на мрежовия протокол.} Приложението кодира и
декодира съобщенията само, срещу публикуваната спецификация~\cite{pg_protocol,
mysql_protocol}. Отпадат всички междинни слоеве, буферът за получаване се притежава от
самата библиотека и декодирането може да работи върху него без копие. Цената е, че
коректността вече е отговорност на проекта, включително частите на протокола, които се
срещат рядко и затова се откриват късно.

\subsection{Критерии за сравнение}

Сравнението се води по критерии, всеки от които е проверим, а не въпрос на вкус: брой
слоеве, през които минава един ред; брой копия на данните между сокета и типизираната
стойност; наличие на естествен асинхронен модел; възможност за кеширане на подготвени
заявления от страна на клиента; обем на кода, който проектът поема да поддържа;
чувствителност към промени във версията на сървъра.

\begin{table}[H]
\centering
\caption{Съпоставка на разгледаните алтернативи по приетите критерии}
\label{tab:alternatives}
\begin{tabular}{L{4.2cm}L{3.0cm}L{3.0cm}L{3.0cm}}
\toprule
\textbf{Критерий} & \textbf{А. ODBC} & \textbf{Б. Доставчикова} & \textbf{В. Пряк протокол} \\
\midrule
Слоеве до реда        & \TODO{?} & \TODO{?} & \TODO{?} \\
Копия на данните      & \TODO{?} & \TODO{?} & \TODO{?} \\
Асинхронен модел      & \TODO{?} & \TODO{?} & \TODO{?} \\
Кеш на заявленията    & \TODO{?} & \TODO{?} & \TODO{?} \\
Поддържан код         & \TODO{?} & \TODO{?} & \TODO{?} \\
\bottomrule
\end{tabular}
\end{table}

Съпоставката по тези критерии \tabref{tab:alternatives} се попълва след като
реализацията и измерванията приключат, за да не съдържа предположения.

\subsection{Обосновка на избора}

Проектът избира алтернатива В. Основанието не е производителност сама по себе си, а
това, че само тя позволява цената на достъпа да бъде наблюдавана. При А и Б броят
копия и обръщания е скрит зад чужд код и може да се оцени само косвено. При В всяко
копие е ред в собствения изходен текст и следователно може да бъде премахнато или
защитено с аргумент. Дисциплината изисква и сравнение между слоестия модел и прякото
говорене по протокол, което тази постановка прави естествено.

Изборът се приема заедно с цената му. Проектът поема поддръжката на две независими
кодиращи схеми, поема риска от разминаване с бъдещи версии на сървърите и се отказва
от преносимостта към източници на данни извън двата разглеждани. Тези ограничения са
записани изрично, за да бъдат отчетени в заключението, а не преоткрити там.
